% concatenated from Galois_symbols.tex
\documentclass[a4paper,12pt]{amsart}

\usepackage{amssymb,amsmath,amsthm,mathrsfs,xspace, braket}
\usepackage[margin=1.0in]{geometry}
\numberwithin{equation}{section}
\usepackage[colorlinks=true]{hyperref}
\usepackage{aliascnt}
\usepackage{mathtools}
\mathtoolsset{showonlyrefs=true}
\usepackage{enumitem}
\setlist[enumerate]{itemsep=0pt,label=$({\rm\roman*})$,topsep=5pt}
\setlist[itemize]{itemsep=0pt,topsep=5pt,labelindent=\parindent,leftmargin=*}
\setlist[description]{itemsep=0pt,topsep=5pt,leftmargin=*}

\usepackage[all,2cell]{xy}
\objectmargin+{1mm}
\labelmargin+{0.8mm}
\SelectTips{cm}{12}
\setcounter{tocdepth}{2}

%\usepackage{showkeys}



\newtheorem{thm}{Theorem}[section]
\newcommand{\thmautorefname}{Theorem}
\newaliascnt{cor}{thm}
\newtheorem{cor}[cor]{Corollary}
\aliascntresetthe{cor}
\newcommand{\corautorefname}{Corollary}
\newaliascnt{lem}{thm}
\newtheorem{lem}[lem]{Lemma}
\aliascntresetthe{lem}
\newcommand{\lemautorefname}{Lemma}
\newaliascnt{prop}{thm}
\newtheorem{prop}[prop]{Proposition}
\aliascntresetthe{prop}
\newcommand{\propautorefname}{Proposition}
\newaliascnt{conj}{thm}
\newtheorem{conj}[conj]{Conjecture}
\aliascntresetthe{conj}
\newcommand{\conjautorefname}{Conjecture}

\theoremstyle{definition}
\newaliascnt{dfn}{thm}
\newtheorem{dfn}[dfn]{Definition}
\aliascntresetthe{dfn}
\newcommand{\dfnautorefname}{Definition}
\newaliascnt{rem}{thm}
\newtheorem{rem}[rem]{Remark}
\aliascntresetthe{rem}
\newcommand{\remautorefname}{Remark}
\renewcommand{\sectionautorefname}{Section}

% Macros used in this paper
\newcommand{\Char}{\operatorname{char}}
\newcommand{\Cor}{\operatorname{Cor}}
\newcommand{\Gm}{\mathbb{G}_{m}}
\DeclareMathOperator{\Jac}{Jac}
\newcommand{\Z}{\mathbb{Z}}
\newcommand{\et}{\mathrm{\acute et}}
\newcommand{\HM}{H_{\mathrm M}}
\newcommand{\Het}{H_{\mathrm{\acute et}}}
\newcommand{\cd}{\operatorname{cd}}
\newcommand{\Mtensor}{\mathbin{\overset{\mathrm M}{\otimes}}}

\title{Galois Symbols for a Jacobian and Multiplicative Groups}


\author[T. Hiranouchi]{Toshiro Hiranouchi}\address[T. Hiranouchi]{
Department of Basic Sciences, Graduate School of Engineering, 
Kyushu Institute of Technology, 
1-1 Sensui-cho, Tobata-ku, Kitakyushu-shi, 
Fukuoka, 804-8550 JAPAN}
\email{hira@mns.kyutech.ac.jp}


\author[R. Sugiyama]{Rin Sugiyama} \address[R. Sugiyama]{
Department of Mathematics, Physics and Computer Science, 
Japan Women's University, 2-8-1 Mejirodai, Bunkyo-ku, Tokyo, 112-8681 JAPAN
}\email{sugiyamar@fc.jwu.ac.jp}


\keywords{MSC2020: 19E15, 14C25, 14F20, 11G10; Somekawa $K$-groups, Galois symbols, Jacobian varieties, higher Chow groups, motivic cohomology}

\begin{document}
\date{July 13, 2026}

\begin{abstract}
Let $C$ be a smooth projective geometrically connected curve over a
field $k$ with a $k$-rational point. 
Let $J$ be the Jacobian variety of $C$.
For an integer $r\geq 1$ and a positive integer
$n$ prime to the characteristic of $k$, we prove that the Galois symbol map 
\[
 K(k;J,\Gm,\ldots,\Gm)/n
 \to
 \Het^{r+1}\bigl(k,J[n]\otimes\mu_n^{\otimes r}\bigr)
\]
is injective, where the multiplicative group $\Gm$ occurs $r$ times.  
The proof uses Akhtar's description of higher Chow groups of zero-cycles and the
Beilinson--Lichtenbaum theorem. The case $r=1$ recovers a theorem of
Spiess.  
%We also record a stable-range $l$-divisibility result for Somekawa $K$-groups attached to several Jacobians and copies of $\Gm$.
\end{abstract}
\maketitle

\section{Introduction}

Let $G_1,\ldots,G_s$ be semi-abelian varieties over a field $k$.
Somekawa introduced the $K$-group
$K(k;G_1,\ldots,G_s)$, now called the Somekawa $K$-group,
by generators over finite extensions of $k$ and by projection-formula
and reciprocity relations \cite{Som90}.  If $n$ is prime to
the characteristic $\Char(k)$ of $k$, 
he also defined the Galois symbol map 
\begin{equation}\label{eq:gen}
 S_n(k;G_1,\ldots,G_s)\colon
 K(k;G_1,\ldots,G_s)/n
 \to
 \Het^s\bigl(k,G_1[n]\otimes\cdots\otimes G_s[n]\bigr);
\end{equation}
see \cite[Lemma 1.6]{Som90}. Somekawa conjectured that this map is
injective \cite[Remark 1.7]{Som90}.  The conjecture is not true for
all semi-abelian varieties; see \cite{SY09}.  It is therefore useful
to find natural classes for which the symbol is injective.

Let $C$ be a smooth projective geometrically connected curve over $k$
with a $k$-rational point, and let $J=\Jac_C$ be its Jacobian.  Spiess
proved that
\[
S_n(k;\Gm,J)\colon K(k;\Gm,J)/n
 \to \Het^2(k,\mu_n\otimes J[n])
\]
is injective (\cite[Appendix, Theorem~6.1]{Yam05}).  His proof uses
the Bloch--Ogus theory and the Merkurjev--Suslin theorem.  The purpose of this note is to extend this result to an arbitrary
number of copies of $\Gm$.

\begin{thm}\label{thm:main}
Let $k$ be a field, and let $C$ be a projective smooth geometrically connected curve over $k$ with a $k$-rational point.  Let $J$ be the
Jacobian variety of $C$.  Let $r\geq 1$, and let $n$ be a positive
integer prime to the characteristic $\Char(k)$ of $k$.  Then the Galois symbol map
\begin{equation}\label{eq:main}
 S_n^{J,r}\colon
 K(k;J,\underbrace{\Gm,\ldots,\Gm}_{r})/n
 \to
 \Het^{r+1}\bigl(k,J[n]\otimes\mu_n^{\otimes r}\bigr)
\end{equation}
is injective.
\end{thm}

The permutation of factors sends
$\set{z,a_1,\ldots,a_r}_{E/k}$ to
$\set{a_1,\ldots,a_r,z}_{E/k}$.  After the corresponding permutation
of the coefficient modules, the two Galois symbols differ by the sign
$(-1)^r$.  Hence the theorem is equivalent to the injectivity of
\[
 K(k;\underbrace{\Gm,\ldots,\Gm}_{r},J)/n
 \to 
 \Het^{r+1}\bigl(k,\mu_n^{\otimes r}\otimes J[n]\bigr).
\]
For $r=1$, after permuting the two factors,
\autoref{thm:main} gives another proof of Spiess's theorem \cite[Appendix, Theorem 6.1]{Yam05}.  
The proof reduces to two known results.  First, Akhtar proved that
a mixed $K$-group attached to $CH_0(C)$ and $r$ copies of
$\Gm$ is naturally isomorphic to the higher Chow group $CH^{r+1}(C,r)$ 
(\cite[Theorem~6.1]{Akh04}).  A rational point gives a decomposition of
this mixed $K$-group into
\[
 K_r^M(k)\oplus K(k;J,\Gm,\ldots,\Gm).
\]
Second, the Beilinson--Lichtenbaum comparison theorem implies that the
finite coefficient \'etale cycle map in bidegree $(r+2,r+1)$ is
injective (\cite[Corollary~1.2]{GL01}, \cite[Theorem~6.17]{Voe11}). 

\subsection*{Notation} 
Throughout this paper, a curve over a field means an integral
separated scheme of finite type over that field of dimension one.

\subsection*{Acknowledgements} 
The first author was supported by JSPS KAKENHI Grant Number 24K06672.
The second author was supported by JSPS KAKENHI Grant Number 21K03188.
%This work was done while the second author's stay in Kyushu Institute of Technology. He would like to thank Kyushu Institute of Technology for their hospitality and for providing an excellent environment to study during his stay.

\section{Preliminaries}

Let $C$ be a smooth projective geometrically connected curve over $k$ with $P\in C(k)$, and $J = \mathrm{Jac}_C$ be the Jacobian of $C$.

\subsection{Mixed \texorpdfstring{$K$}{K}-groups and higher Chow groups}
For every field extension $E/k$, we write $P_E$ for the base change of $P$.
In \cite{Akh04}, Akhtar defined the mixed $K$-group
\[
 K_r\bigl(k;CH_0(C),\Gm\bigr)
 :=K\bigl(k;CH_0(C),\underbrace{\Gm,\ldots,\Gm}_{r}\bigr).
\]
It is generated by symbols
$\set{z,a_1,\ldots,a_r}_{E/k}$, where $E/k$ is finite,
$z\in CH_0(C_E)$, and $a_i\in E^\times$.  The defining relations are
the projection-formula and reciprocity relations; see the definition of
the mixed $K$-group in \cite[Section~3]{Akh04}. 
We also define 
\[
K_r(k;J,\Gm) :=  K(k;J,\underbrace{\Gm,\ldots,\Gm}_{r}).
\]
Let $n$ be an integer prime to $\Char(k)$.  For a finite extension $E/k$, let
\[
 \delta_E\colon E^\times/(E^\times)^n\to \Het^1(E,\mu_n)
\]
and
\[
 \delta_{J,E}\colon J(E)/nJ(E)\to \Het^1(E,J[n])
\]
be the Kummer connecting homomorphisms.  The Galois symbol map
\[
 S_n^{J,r}\colon K_r(k;J,\Gm)/n
 \to \Het^{r+1}\bigl(k,J[n]\otimes\mu_n^{\otimes r}\bigr)
\]
is induced on symbols by
\begin{equation}\label{eq:GS}
 S_n^{J,r}\bigl(\set{z,a_1,\ldots,a_r}_{E/k}\bigr)
 =\Cor_{E/k}\bigl(
 \delta_{J,E}(z)\cup\delta_E(a_1)\cup\cdots\cup\delta_E(a_r)
 \bigr);
\end{equation}
see \cite[Lemma~1.6]{Som90}.
The point $P$ gives, for every finite extension $E/k$, the decomposition
\[
 CH_0(C_E)=\Z[P_E]\oplus CH_0(C_E)^0
          =\Z[P_E]\oplus J(E).
\]
By \cite[Corollary~5.4.5]{Akh00} and
\cite[Proposition~1.5]{Som90}, it induces a natural decomposition
\begin{equation}\label{eq:split}
 K_r\bigl(k;CH_0(C),\Gm\bigr)
 \simeq
 K_r^M(k)\oplus
 K_r(k;J,\Gm).
\end{equation}
We use cohomological notation for higher Chow groups.  By the
Nesterenko--Suslin--Totaro theorem, for every field $E$ there is a
natural isomorphism $K_r^M(E)\simeq CH^r(E,r)$ (\cite{NS90,Tot92}).  
We use this isomorphism to identify these two
groups.  Thus the notation $\set{a_1,\ldots,a_r}$ denotes both the
Milnor symbol and the corresponding class in $CH^r(E,r)$.  Since $C$
has dimension one, the group of zero-cycles in degree $r$ is
$CH^{r+1}(C,r)$.

\begin{thm}[{\cite[Theorem~6.1]{Akh04}}]\label{thm:Akh}
For every $r\geq0$, there is a natural isomorphism
\begin{equation}\label{eq:Akh}
 \varphi_r\colon
 CH^{r+1}(C,r)
 \overset{\sim}{\to}
 K_r\bigl(k;CH_0(C),\Gm\bigr).
\end{equation}
\end{thm}

We describe the inverse map of $\varphi_r$ in \eqref{eq:Akh}.
Let $E/k$ be finite and let $p_E\colon C_E\to C$ be the projection.
Via the canonical identification
\[
 E^\times\simeq CH^1(E,1),
\]
we regard each $a_i\in E^\times$ as its pull-back to
$CH^1(C_E,1)$ along the structure morphism of $C_E$.  The inverse map
$\varphi_r^{-1}$ is characterized by
\begin{equation}\label{eq:Akh-sym}
 \varphi_r^{-1}\bigl(\set{z,a_1,\ldots,a_r}_{E/k}\bigr)
 = (p_E)_*\bigl(z\cap a_1\cap\cdots\cap a_r\bigr)
\end{equation}
for $z\in CH_0(C_E)=CH^1(C_E)$ and $a_i\in E^\times$.  Here $\cap$
denotes the intersection product on higher Chow groups
\[
 \cap\colon CH^p(C_E,m)\times CH^q(C_E,n)
 \to CH^{p+q}(C_E,m+n),
\]
and
\[
 (p_E)_*\colon CH^{r+1}(C_E,r)\to CH^{r+1}(C,r)
\]
is the proper push-forward.  Thus
$z\cap a_1\cap\cdots\cap a_r\in CH^{r+1}(C_E,r)$.
%Akhtar writes these intersection products with dots; the notation in
%\eqref{eq:Akh-sym} represents the same class.

Combining \eqref{eq:split} and \autoref{thm:Akh}, we obtain
\begin{equation}\label{eq:CH}
 CH^{r+1}(C,r)
 \simeq
 K_r^M(k)\oplus
 K_r(k;J,\Gm).
\end{equation}
Let $\iota_{J,r}$ denote the inclusion of the second factor in
\eqref{eq:CH}.  By \eqref{eq:Akh-sym},
\begin{equation}\label{eq:iota}
 \iota_{J,r}\bigl(\set{z,a_1,\ldots,a_r}_{E/k}\bigr)
 = (p_E)_*\bigl(z\cap a_1\cap\cdots\cap a_r\bigr)
\end{equation}
for $z\in J(E)=CH_0(C_E)^0$.
We will see the decomposition \eqref{eq:CH} via a decomposition of the Chow motive of $C$ (\autoref{lem:decomp_CH}).






\subsection{Chow motives and the action of correspondences}
We briefly recall effective Chow motives.
An effective Chow motive is a pair $M=(X,\pi)$ where $X$ is a smooth projective equidimensional variety over $k$ of $\dim X=d$ and $\pi$ is a projector in the ring $CH^{d}(X\times X)$ of correspondences on $X$.
We write $h(X)$ for the Chow motive $(X,1_X)$ where $1_X$ is the graph of the identity $\mathrm{id}_X$. 


A correspondence $\pi\in CH^d(X\times X)$ induces an endomorphism $\pi_\ast$ of $CH^i(X,j)$ as follows.
For any $z\in CH^i(X,j)$, 
$$
\pi_\ast(z):={p_2}_\ast(\pi \cap p_1^\ast(z))\in CH^i(X,j)
$$
where $p_i:X\times X\to X$ is the $i$-th projection.
Note that ${1_X}_\ast$ is the identity.
For an effective Chow motive $M=(X,\pi)$, we define 
$
CH^i(M,j):=\pi_\ast CH^i(X,j)$.
Then we have a decomposition of $CH^i(X,j)$;
$$
CH^i(X,j)=CH^i(M,j)\oplus CH^i(M',j)
$$
where $M'=(X,1_X-\pi)$.


A $k$-rational point $P\in X(k)$ gives the following decomposition of $h(X)$;
$$
h(X)=h^0(X) \oplus h^{[1,2d-1]}(X) \oplus h^{2d}(X) 
$$
where 
\begin{align*}
&h^0(X)=(X, [P\times X]),\\
&h^{[1,2d-1]}(X)=(X, 1_X-[P\times X]-[X\times P])\\
&h^{2d}(X)=(X, [X\times P]).
\end{align*}
There are isomorphisms of Chow motives
$$
h^0(X)\simeq (\operatorname{Spec}(k), 1_{\operatorname{Spec}(k)}),\ h^{2d}(X)\simeq \mathbb{L}^{d}
$$
where $\mathbb{L}$ is the Lefschetz motive.
If $\dim X=1$, we write $h^1(X)$ for $h^{[1,2d-1]}(X)$, and in this case we have the following decomposition of $CH^i(X,j)$;
\begin{align}
    CH^i(X,j)=CH^i(h^0(X),j)\oplus CH^i(h^1(X),j)\oplus CH^i(h^2(X),j).
\end{align}


\begin{lem}\label{lem:decomp_CH}
	For any $r\ge 0$, the isomorphism \eqref{eq:Akh} induces isomorphisms 
    \begin{align}
    &CH^{r+1}(h^0(C),r)=0,\\
    &CH^{r+1}(h^1(C),r)\simeq K_r(k;J,\Gm),\ \text{and}\\
    &CH^{r+1}(h^2(C),r)\simeq K_r^M(k).
    \end{align}
\end{lem}
\begin{proof}
%Since $h^0(C)\simeq (\operatorname{Spec}(k),1_k)$ and $h^2(C)\simeq\mathbb L$,
%we have
%\[
% CH^{r+1}(h^0(C),r)
% \simeq CH^{r+1}(k,r)=0
%\]
%and
%\[
% CH^{r+1}(h^2(C),r)
% \simeq CH^r(k,r)
% \simeq K_r^M(k).
%\]
%We identify the latter group with the Milnor $K$-theory summand in
%\eqref{eq:CH}. 
%
%Let $f\colon C\to\operatorname{Spec}(k)$  and $i\colon P\hookrightarrow C$ 
%be the structure morphism and the closed immersion, respectively.
%Let
%$\gamma\colon C\to C\times C$ given by $x\mapsto(x,P)$.
%Then
%\[
% \gamma_*[C]=[C\times P],\qquad
% p_1\circ\gamma=\operatorname{id}_C,\qquad
% p_2\circ\gamma=i\circ f.
%\]
%Hence, for $z\in CH^a(C,b)$, the projection formula gives
%\[
%\begin{aligned}
% (\pi_2)_*(z)
% &=(p_2)_*\bigl([C\times P]\cap p_1^*z\bigr)\\
% &=(p_2)_*\gamma_*\bigl(\gamma^*p_1^*z\bigr)\\
% &=(p_2\circ\gamma)_*(z)
% =i_*f_*(z).
%\end{aligned}
%\]
%Thus $(\pi_2)_*=i_*f_*$.
%
%Let $E/k$ be finite and let
%\[
% w=(p_E)_*
% \bigl(z\cap a_1\cap\cdots\cap a_r\bigr)
% \in CH^{r+1}(C,r),
%\]
%where $z\in CH_0(C_E)$ and $a_1,\ldots,a_r\in E^\times$.
%By the projection formula and the compatibility of proper
%push-forward with the Nesterenko--Suslin--Totaro isomorphism, we have
%\[
% f_*(w)
% =
% N_{E/k}\bigl(
% \deg_E(z)\{a_1,\ldots,a_r\}
% \bigr)
% \in K_r^M(k).
%\]
%Consequently,
%\[
% (\pi_2)_*(w)
% =
% i_*N_{E/k}\bigl(
% \deg_E(z)\{a_1,\ldots,a_r\}
% \bigr).
%\]
%Under Akhtar's isomorphism \eqref{eq:Akh}, this is the projection onto
%the summand $K_r^M(k)$ in \eqref{eq:CH}. Therefore, the kernel of
%$(\pi_2)_*$ corresponds to the summand $K_r(k;J,\Gm)$.
%
%Since the $h^0(C)$-part of $CH^{r+1}(C,r)$ is zero, the kernel of
%$(\pi_2)_*$ is precisely its $h^1(C)$-part. Hence
%\[
% CH^{r+1}(h^1(C),r)
% \simeq K_r(k;J,\Gm).
%\]
%
    Put $\alpha:=[P\times C], {}^t\alpha:=[C\times P]$.
    We first compute $\alpha_\ast, {}^t\alpha_\ast$ on $CH^1(C)$.
    For $z\in CH^1(C)$, we have
    $$
    \alpha\cap (z\times C)=0,\ {}^t\alpha \cap (z\times C)=z\times P
    $$
    which shows
    $$
    \alpha_\ast(z)=0,\ {}^t\alpha_\ast(z)=\deg(z)[P]\in \Z[P].
    $$
    This shows that the kernel of ${}^t\alpha_\ast$ on $CH^1(C)$ is equal to the kernel of the degree map
    $$
    \deg:CH^1(C) \to \Z
    $$ and proves the assertion in the case $r=0$.
    

    By \eqref{eq:Akh} and \eqref{eq:Akh-sym}, we take an element $w$ of $CH^{r+1}(C,r)$ of the following form
    $$
    w:=(p_E)_\ast\bigl(z\cap a_1\cap\cdots\cap a_r\bigr).
    $$
    By the base change formula for algebraic cycles with respect to proper push-forwards and flat pull-backs, it suffices to compute $\alpha_\ast(w)$ in case $E=k$.
    Then, we have
    \begin{align}
        \alpha_\ast(w)
        &={p_2}_\ast\bigl(\alpha\cap p_1^\ast(z\cap a_1\cap\cdots\cap a_r)\bigr)\\
        &={p_2}_\ast\bigl((\alpha\cap p_1^\ast z)\cap p_1^\ast(a_1\cap\cdots\cap a_r)\bigr)=0,
    \end{align}
    and
    \begin{align}
        {}^t\alpha_\ast(w)
        &={p_2}_\ast\bigl({}^t\alpha\cap p_1^\ast(z\cap a_1\cap\cdots\cap a_r)\bigr)\\
        &={p_2}_\ast\bigl(({}^t\alpha\cap p_1^\ast z)\cap p_1^\ast(a_1\cap\cdots\cap a_r)\bigr)\\
        &={p_2}_\ast\bigl([z\times P]\cap p_1^\ast(a_1\cap\cdots\cap a_r)\bigr)\\
        &=\deg(z)\bigl([P]\cap a_1\cap\cdots\cap a_r\bigr)\\
        &=\deg(z)\{a_1, \dots, a_r\}\in K_r^M(k).
    \end{align}
    Here in the last equality, we identify $CH^r(P,r)$ with $K^M_r(k)$.
    This shows that the kernel of ${}^t\alpha_\ast$ on $CH^{r+1}(C,r)$ is equal to the kernel of the map
    $$
    CH^{r+1}(C,r)\to K_r^M(k).
    $$
    Thus, the isomorphism \eqref{eq:Akh} induces
    $$
    CH^{r+1}(h^0(C),r)=0,\ CH^{r+1}(h^1(C),r)\simeq K_r(k;J,\Gm),\ CH^{r+1}(h^2(C),r)\simeq K_r^M(k).
    $$
\end{proof}
By \autoref{lem:decomp_CH}, the inclusion $\iota_{J,r}$ in \eqref{eq:iota} coincides with 
\begin{align}\label{eq:iota2}
    \iota_{J,r}: K_r(k;J,\Gm)\simeq CH^{r+1}(h^1(C),r)\hookrightarrow CH^{r+1}(C,r).
\end{align}


\subsection{\'Etale cycle map}
For a smooth scheme $X$, we use the identification
\[
 CH^q(X,2q-p)=\HM^p(X,\Z(q));
\]
see \cite{Blo86} and \cite[Theorem~1]{Voe02}.  Let
\begin{equation}\label{eq:rho}
 \rho_{X,n}^{i,j}\colon
 CH^{i}(X,j)/n
 \to
 \Het^{2i-j}\bigl(X,\mu_n^{\otimes i}\bigr)
\end{equation}
be the finite coefficient \'etale cycle map.  For $i=j+1$, its
injectivity follows from \cite[Corollary~1.2]{GL01} (see also \cite[Theorem~6.17]{Voe11}).  
Under the Nesterenko--Suslin--Totaro isomorphism
$CH^r(E,r)\simeq K_r^M(E)$, the finite coefficient \'etale cycle map is the
norm residue isomorphism
\[
 K_r^M(E)/n\overset{\sim}{\to}H^r(E,\mu_n^{\otimes r}),
 \qquad
 \set{a_1,\ldots,a_r}\mapsto
 \delta_E(a_1)\cup\cdots\cup\delta_E(a_r).
\]
%For $n=l$ this follows from \cite[Theorem~6.16]{Voe11}; prime powers
%and general $n$ follow from the coefficient triangles and the Chinese
%remainder theorem as above.
Moreover, under the natural isomorphisms $CH^1(X)\simeq \Het^1(X, \Gm)$ and $CH^1(X,1)\simeq \Gamma(X,\Gm)$, the cycle maps $\rho^{1,0}_{X,n}, \rho^{1,1}_{X,n}$ coincide with the Kummer connecting homomorphisms
\begin{align}\label{eq:1-cyc}
    &\Het^1(X, \Gm)/n\to \Het^2(X,\mu_n),\\
    &\Gamma(X,\Gm)/n\to \Het^1(X,\mu_n).
\end{align}



\begin{lem}\label{lem:cycl}
    Let $X$ be a smooth projective variety over $k$ of dimension $d$.
    For any $\alpha\in CH^d(X\times X)$, put $\rho(\alpha):=\rho_{X\times X,n}^{d,0}(\alpha)$ in $\Het^{2d}(X\times X,\mu_n^{\otimes d})$. Then the following diagram commutes:
    $$
    \xymatrix{
    CH^i(X,j)/n \ar[r]^-{\rho^{i,j}_{X,n}} \ar[d]^{\alpha_*} & H_\et^{2i-j}(X,\mu_n^{\otimes i})\ar[d]^{\rho(\alpha)_*}\\
    CH^i(X,j)/n \ar[r]^-{\rho^{i,j}_{X,n}}  & H_\et^{2i-j}(X,\mu_n^{\otimes i})
    }
    $$
\end{lem}
\begin{proof}
    Let $p_i: X\times X\to X$ be the $i$-th projection.
    Take $z\in CH^i(X,j)$.
    By the compatibility of the cycle map with proper push-forward and flat pull-back, we have
\begin{align*}
 \rho^{i,j}_{X,n}\bigl(\alpha_\ast(z)\bigr)
 &=
 \rho^{i,j}_{X,n}
 \bigl({p_2}_\ast(\alpha\cap p_1^\ast z)\bigr)\\
 &=
 {p_2}_\ast
 \left(
 \rho^{d+i,j}_{X\times X,n}
 \bigl(\alpha\cap p_1^\ast z\bigr)
 \right)\\
 &=
 {p_2}_\ast
 \left(
 \rho(\alpha)\cup
 \rho^{i,j}_{X\times X,n}(p_1^\ast z)
 \right)\\
 &=
 {p_2}_\ast
 \left(
 \rho(\alpha)\cup
 p_1^\ast\bigl(\rho^{i,j}_{X,n}(z)\bigr)
 \right)\\
 &=
 \rho(\alpha)_\ast
 \bigl(\rho^{i,j}_{X,n}(z)\bigr).
\end{align*}
\end{proof}

For an effective Chow motive $M=(X,\pi)$, we define $H_\et^a(M,\mu_n^{\otimes b}):=\rho(\pi)_\ast H_\et^a(X,\mu_n^{\otimes b})$.
By \autoref{lem:cycl}, we have a \'etale cycle map for $M$;
$$
\rho^{i,j}_{M,n}:CH^i(M,j)/n\to H_\et^{2i-j}(M,\mu_n^{\otimes i}).
$$

\begin{lem}[\text{\cite[Lemma 2.5]{Yam05}}]\label{lem:decomp_Het}
%    Let $C$ be a smooth projective geometrically connected curve over $k$.
    For $t\ge 1$, a $k$-rational point $P$ in $C$ gives
    \begin{align}
    &H^{s}_\et(h^0(C),\mu_n^{\otimes t})
    \simeq \Het^{s}(k,\mu_n^{\otimes t}),\\
    &H^s_\et(h^1(C),\mu_n^{\otimes t})
    \simeq H_\et^{s-1}(k,J[n]\otimes\mu_n^{\otimes (t-1)}),\\
    &H^{s}_\et(h^2(C),\mu_n^{\otimes t})
    \simeq H_\et^{s-2}(k,\mu_n^{\otimes (t-1)}).
    \end{align}
\end{lem}
\begin{proof}
    Put
    $\alpha:=[P\times C]$,
    ${}^t\alpha:=[C\times P]$, and 
    $\pi:=1_C-\alpha-{}^t\alpha$.
    Thus
    \[
    h^0(C)=(C,\alpha),\qquad
    h^1(C)=(C,\pi),\qquad
    h^2(C)=(C,{}^t\alpha).
    \]
    Since $
    h^0(C)\simeq (\operatorname{Spec}(k), 1_k), h^2(C)\simeq \mathbb{L},
    $
    the statement is clear for $h^0(C)$ and $h^2(C)$.
    Indeed, for $\rho(\alpha)=\rho^{1,0}_{C\times C,n}([P\times C])$, we have
    $$
    \rho(\alpha)_\ast: \Het^s(C,\mu_n^{\otimes t})\xrightarrow{i^\ast} \Het^s(k,\mu_n^{\otimes t})\xrightarrow{f^\ast} \Het^s(C,\mu_n^{\otimes t})
    $$
    and
    $$
    \rho({}^t\alpha)_\ast: \Het^s(C,\mu_n^{\otimes t})\xrightarrow{f_\ast} \Het^{s-2}(k,\mu_n^{\otimes (t-1)})\xrightarrow{i_\ast} \Het^s(C,\mu_n^{\otimes t}),
    $$
    where $i:P\hookrightarrow C$ is the closed immersion and $f:C\to \operatorname{Spec}(k)$ is the structure map.
    
    For $h^1(C)$, we use the Hochschild--Serre spectral sequence
    \[
    E_2^{a,b}
    =\Het^a\bigl(k, \Het^b(\overline C,
     \mu_n^{\otimes t})\bigr)
 \Longrightarrow
 \Het^{a+b}\bigl(C,\mu_n^{\otimes t}\bigr).
\]
See \cite[Theorem~12.7 and Example~12.8]{Mil13} and
\cite[Chapter~II, Section~4, Theorem~2.4.1]{NSW08}.  This spectral
sequence is multiplicative; on its $E_2$-page the product is induced
by the cup products in continuous Galois cohomology and geometric
\'etale cohomology.  See \cite[pp.~118--119]{HS53} and
\cite[Chapter~I, Section~4]{NSW08}.

The Kummer sequence and the trace isomorphism give
\[
 \Het^0(\overline C, \mu_n^{\otimes t}) \simeq \mu_n^{\otimes t},\ \Het^1(\overline C,\mu_n^{\otimes t})
 \simeq J[n]\otimes\mu_n^{\otimes (t-1)},\ \mbox{and} 
 \
 \Het^2(\overline C,\mu_n^{\otimes t})
 \simeq\mu_n^{\otimes (t-1)};
\]
see \cite[Proposition~14.2 and Theorem~24.1]{Mil13} and
\cite[Sections~8.2--8.3]{Fu15}. 

    The actions of $\alpha$ and ${}^t\alpha$ are compatible with the
    Hochschild--Serre spectral sequence. On the $E_2$-page, their
    actions are induced by the actions of the corresponding
    correspondences on $\Het^b(\overline C,\mu_n^{\otimes t})$.
    After base change to $\overline k$, the same formulas as above show
    that $\rho(\alpha)_\ast$ is the identity on
    $\Het^0(\overline C,\mu_n^{\otimes t})$ 
    and is zero on
    $\Het^1(\overline C,\mu_n^{\otimes t})$ 
    and $\Het^2(\overline C,\mu_n^{\otimes t})$.
    Similarly, $\rho({}^t\alpha)_\ast$ is the identity on
    $\Het^2(\overline C,\mu_n^{\otimes t})$
    and is zero on
    $\Het^0(\overline C,\mu_n^{\otimes t})$
    and $\Het^1(\overline C,\mu_n^{\otimes t})$.
    Since
    \[
    \rho(\pi)_\ast
    =
    1-\rho(\alpha)_\ast-\rho({}^t\alpha)_\ast,
    \]
    it follows that $\rho(\pi)_\ast$ acts as zero on
    $\Het^0(\overline C,\mu_n^{\otimes t})$ 
    and $\Het^2(\overline C,\mu_n^{\otimes t})$,
    and as the identity on
    $\Het^1(\overline C,\mu_n^{\otimes t})$.
    Therefore the Hochschild--Serre spectral sequence obtained by
    applying $\rho(\pi)_\ast$ has only the row $b=1$, namely
    \[
    E_2^{a,1}
    =
    \Het^a\bigl(
    k,J[n]\otimes\mu_n^{\otimes(t-1)}
    \bigr).
    \]
    Hence there are no nonzero differentials, and we obtain
    \[
    H^s_\et(h^1(C),\mu_n^{\otimes t})
    =
    \rho(\pi)_\ast
    \Het^s(C,\mu_n^{\otimes t})
    \simeq
    \Het^{s-1}\bigl(
    k,J[n]\otimes\mu_n^{\otimes(t-1)}
    \bigr).
    \]
%***
%Therefore, we have the following exact sequences
%\begin{align}
%    0\to F^1\to \Het^s(C, \mu_n^{\otimes t}) \xrightarrow{f_\ast} \Het^{s-2}(k,\mu_n^{\otimes (t-1)})\to 0,\\
%    0\to \Het^{s}(k,\mu_n^{\otimes t})\xrightarrow{f^\ast} F^1\to \Het^{s-1}(k, J[n]\otimes \mu_n^{\otimes (t-1)})\to 0.
%\end{align}
%The Chow--K\"unneth projectors associated with $P$ split the Hochschild--Serre filtration.
This shows the assertion for $h^1(C)$.
\end{proof}

For $r\ge 0$, let $\iota^\et_{J,r}$ denote the inclusion 
\begin{align}\label{eq:iota_et}
 \iota^\et_{J,r}\colon
 \Het^{r+1}\bigl(k,J[n]\otimes\mu_n^{\otimes r}\bigr)
 \hookrightarrow
 \Het^{r+2}\bigl(C,\mu_n^{\otimes(r+1)}\bigr)
\end{align}
induced by the Hochschild--Serre spectral sequence and the
Chow--K\"unneth splitting, with the total-complex convention of
\cite[Chapter~II, Section~2, Definition~2.2.3 and p.~104]{NSW08}.

\begin{lem}\label{lem:KJ}
For every $z\in J(k)=CH_0(C)^0$, one has
\begin{equation}\label{eq:KJ}
 \rho_{C,n}^{1,0}(z)
 =\iota_{J,0}^{\et}\bigl(\delta_J(z)\bigr)
 \quad\text{in }H_\et^2(C,\mu_n).
\end{equation}
Equivalently, under the identification
\[
 H_\et^2\bigl(h^1(C),\mu_n\bigr)\simeq \Het^1(k,J[n]),
\]
the cycle map
\[
 \rho_{h^1(C),n}^{1,0}\colon J(k)/n
 \to \Het^1(k,J[n])
\]
coincides with the Kummer connecting homomorphism $\delta_J$.
\end{lem}

\begin{proof}
Under the identification $CH^1(C)=\operatorname{Pic}(C)$, the map
$\rho_{C,n}^{1,0}$ is the connecting homomorphism associated with the
Kummer sequence on $C$. Since a degree-zero class belongs to the
$h^1(C)$-part, the compatibility of the Kummer sequence with the
Hochschild--Serre edge morphism gives the commutative diagram
\[
\xymatrix{
 J(k)/n\ar[r]^-{\rho_{h^1(C),n}^{1,0}}\ar[d]_{\delta_J}
 &H_\et^2(h^1(C),\mu_n)\ar[d]^{\simeq}\\
 \Het^1(k,J[n])\ar@{=}[r]&\Het^1(k,J[n]).
}
\]
This proves the equivalent assertion, and \eqref{eq:KJ} follows by
including the $h^1(C)$-part into $H_\et^2(C,\mu_n)$.
\end{proof}

\begin{lem}\label{lem:HS-sign}
Let $f\colon C\to\operatorname{Spec}(k)$ be the structure morphism.
For $\alpha\in \Het^1(k,J[n])$ and $\beta\in \Het^r(k,\mu_n^{\otimes r})$, 
one has
\begin{equation}\label{eq:HS-sign}
 \iota_{J,0}^{\et}(\alpha)\cup f^*\beta
 =(-1)^r\iota_{J,r}^{\et}(\alpha\cup\beta).
\end{equation}
\end{lem}
\begin{proof}
In the standard total complex for the Hochschild--Serre spectral
sequence, the product of classes of bidegrees $(a,b)$ and $(a',b')$
acquires the sign $(-1)^{ba'}$; see
\cite[Chapter II, Section~1]{HS53} and
\cite[Chapter~II, Section~4, proof of Theorem~2.4.1]{NSW08}.
Here $\alpha$ has bidegree $(1,1)$ and $\beta$ has bidegree $(r,0)$.
Thus their product acquires the sign
$(-1)^{1\cdot r}=(-1)^r$.
Under the Chow--K\"unneth splitting, this is precisely
\eqref{eq:HS-sign}.
\end{proof}

%\begin{proof}
%Consider the multiplicative Hochschild--Serre spectral sequence
%\[
% E_2^{a,b}(t)
% =H^a\bigl(k,H^b(\overline C,\mu_n^{\otimes t})\bigr)
% \Longrightarrow
% \Het^{a+b}(C,\mu_n^{\otimes t}).
%\]
%Under the identification
%\[
% H^1(\overline C,\mu_n)\simeq J[n],
%\]
%the class $\alpha$ belongs to
%\[
% E_2^{1,1}(1)
% =H^1\bigl(k,H^1(\overline C,\mu_n)\bigr),
%\]
%whereas $\beta$ belongs to
%\[
% E_2^{r,0}(r)
% =H^r\bigl(k,H^0(\overline C,\mu_n^{\otimes r})\bigr).
%\]
%In the construction of the Hochschild--Serre spectral sequence, the
%differential in the second direction is multiplied by $(-1)^a$ on
%the component of bidegree $(a,b)$; see
%\cite[Chapter~II, Section~4, proof of Theorem~2.4.1,
%pp.~111--112]{NSW08}.  Together with the multiplicative structure of
%the spectral sequence \cite[pp.~118--119]{HS53}, this convention gives
%the Koszul sign $(-1)^{ba'}$ for the product of classes of bidegrees
%$(a,b)$ and $(a',b')$.  Here the relevant bidegrees are $(1,1)$ and
%$(r,0)$, so the sign is $(-1)^{1\cdot r}=(-1)^r$.  Therefore the
%product $\iota_{J,0}^{\et}(\alpha)\cup f^*\beta$ corresponds to
%$(-1)^r\alpha\cup\beta$ in
%\[
% H^{r+1}\bigl(k,J[n]\otimes\mu_n^{\otimes r}\bigr),
%\]
%which proves \eqref{eq:HS-sign}.
%\end{proof}



\section{Proof of the main theorem}


\begin{proof}[Proof of \autoref{thm:main}]
By \autoref{lem:decomp_CH}, the map $\iota_{J,r}$ is the inclusion
of a direct summand. Hence it remains injective after reduction
modulo $n$. From \autoref{lem:decomp_CH} and
\autoref{lem:decomp_Het}, we therefore have a commutative diagram
\[
\xymatrix{
 K_r(k;J,\Gm)/n\ar@{^{(}->}[r]^{\iota_{J,r}}_{\eqref{eq:iota2}}\ar[d]^{\rho_{h^1(C),n}^{r+1,r}}
 &CH^{r+1}(C,r)/n\ar@{^{(}->}[d]^{\rho_{C,n}^{r+1,r}}\\
 \Het^{r+1}\bigl(k,J[n]\otimes\mu_n^{\otimes r}\bigr)
 \ar@{^{(}->}[r]^-{\iota_{J,r}^{\et}}_{\eqref{eq:iota_et}}
 &\Het^{r+2}\bigl(C,\mu_n^{\otimes(r+1)}\bigr).
}
\]
The right vertical map is injective by
\cite[Corollary~1.2]{GL01} (see also
\cite[Theorem~6.17(1)]{Voe11}); hence the left vertical map
\[
 \rho_{h^1(C),n}^{r+1,r}\colon
 K_r(k;J,\Gm)/n
 \longrightarrow \Het^{r+1}\bigl(k,J[n]\otimes\mu_n^{\otimes r}\bigr)
\]
is also injective.

It remains to compare the \'etale cycle map
$\rho_{h^1(C),n}^{r+1,r}$ with the Galois symbol map
$S_n^{J,r}$. We first consider symbols defined over the base field.
Let $f\colon C\to\operatorname{Spec}(k)$ be the structure morphism.
For $z\in J(k)$ and $a_1,\ldots,a_r\in k^\times$, the compatibility
of the cycle map with the intersection product and the cup product gives
\begin{align*}
 \iota_{J,r}^{\et}\circ\rho_{h^1(C),n}^{r+1,r}
 \bigl(\{z,a_1,\ldots,a_r\}_{k/k}\bigr)
 &=\rho_{C,n}^{r+1,r}
 \bigl(z\cap a_1\cap\cdots\cap a_r\bigr)\\
 &=\rho_{C,n}^{1,0}(z)\cup
   \rho_{C,n}^{1,1}(a_1)\cup\cdots\cup
   \rho_{C,n}^{1,1}(a_r)\\
 &\stackrel{\eqref{eq:KJ}}{=}
 \iota_{J,0}^{\et}\bigl(\delta_J(z)\bigr)\cup
 f^*\delta(a_1)\cup\cdots\cup f^*\delta(a_r)\\
 &\stackrel{\eqref{eq:HS-sign}}{=}
 (-1)^r\iota_{J,r}^{\et}\bigl(
 \delta_J(z)\cup\delta(a_1)\cup\cdots\cup\delta(a_r)
 \bigr).
\end{align*}
Thus the desired comparison holds for symbols over $k$.

Now let $E/k$ be a finite extension, and let
$p_E\colon C_E\to C$ be the projection. Write
$\iota_{J_E,r}^{\et}$ for the corresponding inclusion over $E$.
The naturality of the Hochschild--Serre spectral sequence and the
compatibility of proper push-forward with trace give a commutative
diagram
\begin{equation}\label{eq:trace}
\vcenter{
\xymatrix{
 \Het^{r+1}\bigl(E,J[n]\otimes\mu_n^{\otimes r}\bigr)
 \ar[r]^-{\iota_{J_E,r}^{\et}}\ar[d]_{\Cor_{E/k}}
 &\Het^{r+2}\bigl(C_E,\mu_n^{\otimes(r+1)}\bigr)
 \ar[d]^{(p_E)_*}\\
 \Het^{r+1}\bigl(k,J[n]\otimes\mu_n^{\otimes r}\bigr)
 \ar[r]^-{\iota_{J,r}^{\et}}
 &\Het^{r+2}\bigl(C,\mu_n^{\otimes(r+1)}\bigr).
}}
\end{equation}
See \cite[Expos\'e~IX, Section~5]{SGA4} and
\cite[Section~8.2]{Fu15}.
Applying the comparison over $E$ and using the compatibility of the
cycle map with proper push-forward, we obtain, for
$z\in J(E)$ and $a_1,\ldots,a_r\in E^\times$,
\begin{align*}
 &\iota_{J,r}^{\et}\circ\rho_{h^1(C),n}^{r+1,r}
 \bigl(\{z,a_1,\ldots,a_r\}_{E/k}\bigr)\\
 &\quad=
 \rho_{C,n}^{r+1,r}\Bigl(
 (p_E)_*\bigl(z\cap a_1\cap\cdots\cap a_r\bigr)
 \Bigr)\\
 &\quad=
 (p_E)_*\rho_{C_E,n}^{r+1,r}
 \bigl(z\cap a_1\cap\cdots\cap a_r\bigr)\\
 &\quad=
 (-1)^r(p_E)_*\iota_{J_E,r}^{\et}\bigl(
 \delta_{J,E}(z)\cup\delta_E(a_1)\cup\cdots\cup\delta_E(a_r)
 \bigr)\\
 &\quad\stackrel{\eqref{eq:trace}}{=}
 (-1)^r\iota_{J,r}^{\et}\Cor_{E/k}\bigl(
 \delta_{J,E}(z)\cup\delta_E(a_1)\cup\cdots\cup\delta_E(a_r)
 \bigr)\\
 &\quad\stackrel{\eqref{eq:GS}}{=}
 (-1)^r\iota_{J,r}^{\et}\circ S_n^{J,r}
 \bigl(\{z,a_1,\ldots,a_r\}_{E/k}\bigr).
\end{align*}
Since such symbols generate $K_r(k;J,\Gm)$ and
$\iota_{J,r}^{\et}$ is injective, we obtain
\[
 \rho_{h^1(C),n}^{r+1,r}=(-1)^rS_n^{J,r}.
\]
Since $\rho_{h^1(C),n}^{r+1,r}$ is injective, the Galois symbol map
$S_n^{J,r}$ is injective.
\end{proof}


% \begin{proof}[Proof of \autoref{thm:main}]
% For a finite extension $E/k$, choose a separable closure $\overline E$
% and put $\overline C_E=C_E\times_E\overline E$.  For $t=0,r$, consider
% the Hochschild--Serre spectral sequence
% \[
%  E_2^{a,b}(E,t)
%  =H^a\bigl(E,H^b(\overline C_E,
%      \mu_n^{\otimes(t+1)})\bigr)
%  \Longrightarrow
%  \Het^{a+b}\bigl(C_E,\mu_n^{\otimes(t+1)}\bigr).
% \]
% See \cite[Theorem~12.7 and Example~12.8]{Mil13} and
% \cite[Chapter~II, Section~4, Theorem~2.4.1]{NSW08}.  This spectral
% sequence is multiplicative; on its $E_2$-page the product is induced
% by the cup products in continuous Galois cohomology and geometric
% \'etale cohomology.  See \cite[pp.~118--119]{HS53} and
% \cite[Chapter~I, Section~4]{NSW08}.

% The Kummer sequence and the trace isomorphism give
% \[
%  H^1(\overline C_E,\mu_n^{\otimes(t+1)})
%  \simeq J[n]\otimes\mu_n^{\otimes t},
%  \qquad
%  H^2(\overline C_E,\mu_n^{\otimes(t+1)})
%  \simeq\mu_n^{\otimes t};
% \]
% see \cite[Proposition~14.2 and Theorem~24.1]{Mil13} and
% \cite[Sections~8.2--8.3]{Fu15}.  The point $P_E$ yields the natural
% split decomposition
% \begin{equation}\label{eq:HS}
%  \Het^{t+2}\bigl(C_E,\mu_n^{\otimes(t+1)}\bigr)
%  \simeq
%  H^{t+2}\bigl(E,\mu_n^{\otimes(t+1)}\bigr)
%  \oplus H^{t+1}\bigl(E,J[n]\otimes\mu_n^{\otimes t}\bigr)
%  \oplus H^t\bigl(E,\mu_n^{\otimes t}\bigr);
% \end{equation}
% see \cite[Lemma~2.5]{Yam05}.  Let
% \[
%  \iota_{J,E}^{\et,0}\colon
%  H^1(E,J[n])
%  \to
%  \Het^2(C_E,\mu_n)
% \]
% and
% \[
%  \iota_{J,E}^{\et,r}\colon
%  H^{r+1}\bigl(E,J[n]\otimes\mu_n^{\otimes r}\bigr)
%  \to
%  \Het^{r+2}\bigl(C_E,\mu_n^{\otimes(r+1)}\bigr)
% \]
% be the inclusions given by \eqref{eq:HS}.  For
% $E=k$, we write $\iota_J^{\et,0}$ and $\iota_J^{\et,r}$.

% It is enough to prove the commutativity of
% \[
% \xymatrix{
%  CH^{r+1}(C,r)/n\ar@{^{(}->}[r]^{\rho_{C,n}^r}
%  & \Het^{r+2}\bigl(C,\mu_n^{\otimes(r+1)}\bigr)\\
%  K_r(k;J,\Gm)/n\ar@{^{(}->}[u]^{\iota_J}\ar[r]^{S_n^{J,r}}
%  & H^{r+1}\bigl(k,J[n]\otimes\mu_n^{\otimes r}\bigr)
%    \ar@{^{(}->}[u]_{\iota_J^{\et,r}} .
% }
% \]
% For a finite extension $E/k$, let
% $q_E\colon C_E\to\operatorname{Spec}(E)$ be the structure morphism.
% We use the finite coefficient cycle maps
% \[
%  \rho_{C_E,n}^{0}\colon
%  CH^1(C_E)/n\to \Het^2(C_E,\mu_n)
% \]
% and
% \[
%  \rho_{C_E,n}^{1,1}\colon
%  CH^1(C_E,1)/n\to \Het^1(C_E,\mu_n).
% \]
% They satisfy
% \begin{align}
%  \rho_{C_E,n}^r\bigl(z\cap a_1\cap\cdots\cap a_r\bigr)
%  &=\rho_{C_E,n}^0(z)\cup
%    \rho_{C_E,n}^{1,1}(a_1)\cup\cdots\cup
%    \rho_{C_E,n}^{1,1}(a_r),
%  \label{eq:prod}\\
%  \rho_{C,n}^r\bigl((p_E)_*u\bigr)
%  &=(p_E)_*\rho_{C_E,n}^r(u),
%  \qquad u\in CH^{r+1}(C_E,r)/n.
%  \label{eq:push}
% \end{align}
% The first equality is the compatibility of the cycle map with the
% intersection product and the cup product, as used in
% \cite[proof of Proposition~2.4]{Yam05}.  The second follows from the
% compatibility of the motivic-to-\'etale comparison morphism with
% transfers; see \cite[Section~3.7, especially (3.21)]{GL01}.  On the
% \'etale side, the corresponding push-forward is the trace map; see
% \cite[Expos\'e~IX, Section~5]{SGA4} and
% \cite[Section~8.2]{Fu15}.

% For $a\in E^\times=CH^1(E,1)$, one has
% \begin{equation}\label{eq:KG}
%  \rho_{C_E,n}^{1,1}(q_E^*a)
%  =q_E^*\delta_E(a).
% \end{equation}

% \noindent 
% *****Proof of (3.4)****
% Indeed, 
% let $X$ be a smooth scheme on which $n$ is invertible, and let
% $\epsilon\colon X_{\et}\to X_{\mathrm{Zar}}$
% be the canonical morphism of sites.  The finite coefficient \'etale
% cycle map is induced by
% \[
%  \Z/n(1)\to\tau_{\leq1}R\epsilon_*\mu_n;
% \]
% see \cite[Section~1, (1.2)]{GL01}.  Since
% $\Z(1)\simeq\Gm[-1]$,
% the coefficient triangle for $\Z(1)$ is identified on the \'etale
% site with the triangle associated with the Kummer sequence
% \[
%  1\to\mu_n\to\Gm\xrightarrow{\,n\,}\Gm\to1.
% \]
% Hence, under
% $CH^1(X,1)=\Gamma(X,\mathcal O_X^\times)$,
% the map
% \[
%  \rho_{X,n}^{1,1}\colon
%  CH^1(X,1)/n\to\Het^1(X,\mu_n)
% \]
% coincides with the Kummer connecting homomorphism.
% Applying this to
% $q_E\colon C_E\to\operatorname{Spec}(E)$ and using functoriality with
% respect to pull-back, we obtain \eqref{eq:KG}.

% \noindent 
% **** End of the proof of (3.4) ****


% Let $z\in J(E)=CH_0(C_E)^0$.  Then
% \begin{equation}\label{eq:KJ}
%  \rho_{C_E,n}^0(z)
%  =\iota_{J,E}^{\et,0}\bigl(\delta_{J,E}(z)\bigr).
% \end{equation}

% \noindent 
% **** Proof of (3.5) **** 

% Compare the $l$-adic formulation in
% \cite[Appendix, Lemma]{Ras95}???

% \noindent 
% **** End of the proof ****

% The multiplicative structure on the Hochschild--Serre spectral
% sequence gives
% \begin{equation}\label{eq:cup}
%  \iota_{J,E}^{\et,r}(\alpha\cup\beta)
%  =(-1)^r\iota_{J,E}^{\et,0}(\alpha)\cup q_E^*\beta
% \end{equation}
% for $\alpha\in H^1(E,J[n])$ and
% $\beta\in H^r(E,\mu_n^{\otimes r})$.  

% \noindent 
% **** Proof of (3.6) **** 


% Here we use the compatibility of the Hochschild--Serre filtration
% with cup products from \cite[pp.~118--119]{HS53}, together with the
% cup-product formalism of continuous Galois cohomology described in
% \cite[Chapter~I, Section~4]{NSW08}, and the splitting
% \eqref{eq:HS}.

% \noindent 
% **** End of the proof ****


% The naturality of the Hochschild--Serre spectral sequence and the
% trace morphism show that proper push-forward along $p_E$ induces
% Galois-cohomological corestriction on the $b=1$ term.  See
% \cite[Expos\'e~IX, Section~5]{SGA4} and
% \cite[Section~8.2]{Fu15}.  Thus the following diagram is commutative:
% \begin{equation}\label{eq:trace}
% \vcenter{
% \xymatrix{
%  H^{r+1}\bigl(E,J[n]\otimes\mu_n^{\otimes r}\bigr)
%  \ar[r]^-{\iota_{J,E}^{\et,r}}\ar[d]_{\Cor_{E/k}}
%  & \Het^{r+2}\bigl(C_E,\mu_n^{\otimes(r+1)}\bigr)
%    \ar[d]^{(p_E)_*}\\
%  H^{r+1}\bigl(k,J[n]\otimes\mu_n^{\otimes r}\bigr)
%  \ar[r]^-{\iota_J^{\et,r}}
%  & \Het^{r+2}\bigl(C,\mu_n^{\otimes(r+1)}\bigr).
% }}
% \end{equation}

% Take
% $x=\set{z,a_1,\ldots,a_r}_{E/k}$ with $z\in J(E)$.  Then
% \begin{align*}
%  \rho_{C,n}^r\bigl(\iota_J(x)\bigr)
%  &\overset{\eqref{eq:iota}}{=}
%  \rho_{C,n}^r\Bigl(
%    (p_E)_*\bigl(z\cap a_1\cap\cdots\cap a_r\bigr)
%    \Bigr)\\
%  &\overset{\eqref{eq:prod},\,
%              \eqref{eq:push}}{=}
%  (p_E)_*\Bigl(
%    \rho_{C_E,n}^0(z)\cup
%    \rho_{C_E,n}^{1,1}(a_1)\cup\cdots\cup
%    \rho_{C_E,n}^{1,1}(a_r)
%    \Bigr)\\
%  &\overset{\eqref{eq:KG},\,
%              \eqref{eq:KJ}}{=}
%  (p_E)_*\Bigl(
%    \iota_{J,E}^{\et,0}\bigl(\delta_{J,E}(z)\bigr)
%    \cup q_E^*\delta_E(a_1)\cup\cdots\cup q_E^*\delta_E(a_r)
%    \Bigr)\\
%  &\overset{\eqref{eq:cup}}{=}
%  (p_E)_*\iota_{J,E}^{\et,r}\Bigl(
%    \delta_{J,E}(z)\cup\delta_E(a_1)\cup\cdots\cup\delta_E(a_r)
%    \Bigr)\\
%  &\overset{\eqref{eq:trace}}{=}
%  \iota_J^{\et,r}\Cor_{E/k}\Bigl(
%    \delta_{J,E}(z)\cup\delta_E(a_1)\cup\cdots\cup\delta_E(a_r)
%    \Bigr)\\
%  &\overset{\eqref{eq:GS}}{=}
%  \iota_J^{\et,r}S_n^{J,r}(x).
% \end{align*}
% Thus the diagram is commutative.  Since $\rho_{C,n}^r$ and
% $\iota_J$ are injective, the Galois symbol $S_n^{J,r}$ is injective.
% \end{proof}


\begin{rem}[\text{\cite[Lemma~3.2]{HS25}, \cite{Hir24}}]\label{prop:div}
We record a stable-range $l$-divisibility result for Somekawa $K$-groups attached to several semi-abelian varieties and copies of $\Gm$.
Let $G_1,\ldots, G_q\ (q\ge 1)$ be semi-abelian varieties over $k$%Jacobian varieties over $k$
, and let
$l\neq\Char(k)$ be a prime.  Assume that
$s=\cd_l(k)<\infty$.  If $r\ge s$, then
$K(k;G_1,\ldots,G_q,
 \underbrace{\Gm,\ldots,\Gm}_{r})$
is $l$-divisible.  Consequently, its quotient modulo $l^m$ is zero
for every $m\geq1$, and the corresponding Galois symbol map modulo
$l^m$ is injective.
\end{rem}

% \begin{proof}
% Let $\Mtensor$ denote the Mackey product.  By the definition of the
% Somekawa $K$-group, there is a surjection
% \[
% \bigl(J_1\Mtensor\cdots\Mtensor J_q\Mtensor
%        \Gm^{\Mtensor r}\bigr)(k)/l\\
% \twoheadrightarrow
%  K(k;J_1,\ldots,J_q,
%  \underbrace{\Gm,\ldots,\Gm}_{r})/l.
% \]
% The Mackey product is associative and right exact.  Hence Kahn's
% theorem, in the form of \cite[Remark~2.4 and Theorem~4.5]{Hir24},
% gives
% \begin{align*}
% \bigl(J_1\Mtensor\cdots\Mtensor J_q\Mtensor  \Gm^{\Mtensor r})(k)/l&\simeq \bigl(J_1\Mtensor\cdots\Mtensor J_q\Mtensor
%        (\Gm^{\Mtensor r})/l\bigr)(k)\\
%  &\simeq
%  \bigl(J_1\Mtensor\cdots\Mtensor J_q\Mtensor
%        (K_r^M/l)\bigr)(k).
% \end{align*}
% For every finite extension $E/k$, one has $\cd_l(E)\leq s$.
% Such a field is a $\mathfrak B_s(l)$-field in the sense of \cite[Chapter~II]{Kat80} 
% and that $K_t^M(E)$ is $l$-divisible for every $t>s$ (\cite[Chapter~II, Section~3.3, Remark~1 and Lemma~10]{Kat80}).
% Since $r>s$, we have $(K_r^M/l)(E)=0$ for every finite extension
% $E/k$.  Therefore the last Mackey product is zero.  It follows that
% \[
%  K(k;J_1,\ldots,J_q,
%  \underbrace{\Gm,\ldots,\Gm}_{r})/l=0.
% \]
% Thus the Somekawa $K$-group is $l$-divisible.  Repeated division by $l$
% gives the vanishing modulo $l^m$ for all $m\geq1$.
% \end{proof}

\begin{thebibliography}{Yam05}

\bibitem[Akh00]{Akh00}
R.~Akhtar, \emph{Milnor {$K$}-theory and zero-cycles on algebraic varieties},
Thesis available at \url{http://calico.mth.muohio.edu/reza/} (2000).

\bibitem[Akh04]{Akh04}
\bysame, \emph{Milnor {$K$}-theory of smooth varieties}, $K$-Theory
\textbf{32} (2004), no.~3, 269--291.

\bibitem[Blo86]{Blo86}
S.~Bloch, \emph{Algebraic cycles and higher {$K$}-theory}, Adv. Math.
\textbf{61} (1986), no.~3, 267--304.

\bibitem[Fu15]{Fu15}
L.~Fu, \emph{\'Etale cohomology theory}, revised edition,
Nankai Tracts in Mathematics, vol.~14, World Scientific Publishing Co.
Pte. Ltd., Hackensack, NJ, 2015.

\bibitem[GL01]{GL01}
T.~Geisser and M.~Levine,
\emph{The Bloch--Kato conjecture and a theorem of Suslin--Voevodsky},
J. Reine Angew. Math. \textbf{530} (2001), 55--103.

\bibitem[Hir24]{Hir24}
T.~Hiranouchi, \emph{An additive variant of the differential symbol maps}, Ann. K-Theory \textbf{9} (2024), no.~3, 499--518.

\bibitem[HS25]{HS25}
T.~Hiranouchi and R.~Sugiyama, \emph{Zero-cycles on varieties over a $\mathfrak{B}_s$-field}, arXiv:2509.15617v3. 


\bibitem[HS53]{HS53}
G.~Hochschild and J.-P.~Serre, \emph{Cohomology of group extensions},
Trans. Amer. Math. Soc. \textbf{74} (1953), 110--134.



% \bibitem[Kat80]{Kat80}
% K.~Kato, \emph{A generalization of local class field theory by using
% {$K$}-groups. {II}}, J. Fac. Sci. Univ. Tokyo Sect. IA Math. \textbf{27}
% (1980), no.~3, 603--683.

\bibitem[Mil13]{Mil13}
J.~S. Milne, \emph{Lectures on \'Etale Cohomology}, version 2.21,
2013, available at \url{https://www.jmilne.org/math/CourseNotes/LEC.pdf}.

\bibitem[NSW08]{NSW08}
J.~Neukirch, A.~Schmidt, and K.~Wingberg,
\emph{Cohomology of number fields}, second edition,
Grundlehren der Mathematischen Wissenschaften, vol.~323,
Springer-Verlag, Berlin, 2008.

\bibitem[NS90]{NS90}
Y.~P. Nesterenko and A.~A. Suslin, \emph{Homology of the general linear group
over a local ring, and {M}ilnor's {$K$}-theory}, Math. USSR-Izv. \textbf{34}
(1990), no.~1, 121--145.

% \bibitem[Ras95]{Ras95}
% W.~Raskind, \emph{Higher {$l$}-adic {A}bel--{J}acobi mappings and
% filtrations on {C}how groups}, Duke Math. J. \textbf{78} (1995), no.~1,
% 33--57.

% \bibitem[RS00]{RS00}
% W.~Raskind and M.~Spiess, \emph{Milnor {$K$}-groups and zero-cycles on
% products of curves over {$p$}-adic fields}, Compositio Math. \textbf{121}
% (2000), no.~1, 1--33.

\bibitem[SGA4]{SGA4}
M.~Artin, A.~Grothendieck, and J.-L.~Verdier,
\emph{Th\'eorie des topos et cohomologie \'etale des sch\'emas
(SGA 4), Tome 2}, Lecture Notes in Mathematics, vol.~270,
Springer-Verlag, Berlin--New York, 1972.

\bibitem[Som90]{Som90}
M.~Somekawa, \emph{On {M}ilnor {$K$}-groups attached to semi-abelian
varieties}, $K$-Theory \textbf{4} (1990), no.~2, 105--119.

\bibitem[SY09]{SY09}
M.~Spiess and T.~Yamazaki, \emph{A counterexample to generalizations of
the {M}ilnor--{B}loch--{K}ato conjecture}, J. K-Theory \textbf{4} (2009),
no.~1, 77--90.

\bibitem[Tot92]{Tot92}
B.~Totaro, \emph{Milnor {$K$}-theory is the simplest part of algebraic
{$K$}-theory}, $K$-Theory \textbf{6} (1992), no.~2, 177--189.

\bibitem[Voe02]{Voe02}
V.~Voevodsky, \emph{Motivic cohomology groups are isomorphic to higher
{C}how groups in any characteristic}, Int. Math. Res. Not. (2002), no.~7,
351--355.

\bibitem[Voe11]{Voe11}
\bysame, \emph{On motivic cohomology with
{$\mathbb Z/l$}-coefficients}, Ann. of Math. (2) \textbf{174} (2011), no.~1,
401--438.

\bibitem[Yam05]{Yam05}
T.~Yamazaki, \emph{On {C}how and {B}rauer groups of a product of {M}umford
curves}, Math. Ann. \textbf{333} (2005), 549--567.

\end{thebibliography}

\end{document}
